\chapter{The architecture, in three sentences}
\label{ch:summary}

\begin{aside}
You have just spent thirty chapters on Part II. Before turning
the page to Part III, it is worth distilling everything to its
shortest form. If you remember only three things, remember
these.
\end{aside}

\section{Sentence 1: the data flow}

\begin{quote}
A \maya{} program is four pure functions
(\code{init}, \code{update}, \code{view}, \code{subscribe}) over
two sum types (\code{Msg}, \code{Element}) plus two value-typed
effect languages (\code{Cmd}, \code{Sub}).
\end{quote}

Everything in Part II descended from that sentence. The Elm
loop, the DSL, elements, layout, concepts, expected, the
pipeline, style, builders, text, focus, scroll, error
boundaries, signals, the runtime, widgets, performance,
patterns --- all of them are either inside one of those types
or transformations between them.

\subsection*{Why the four functions}

Each one separates a concern that becomes a bug when entangled:

\begin{itemize}[leftmargin=*]
  \item \code{init}: how does the program start? Separated so
    boot-time state has its own home.
  \item \code{update}: how does state evolve? Separated so
    business logic is testable in isolation.
  \item \code{view}: how does state become pixels? Separated so
    the renderer is a pure function of state.
  \item \code{subscribe}: how does the program listen?
    Separated so external events are declarative.
\end{itemize}

Combine any two and you lose the property. Keeping them apart
is the discipline.

\subsection*{Why two sum types}

\code{Msg} is the alphabet of state transitions. Each
alternative is one kind of thing that can happen. The compiler
checks exhaustiveness; new alternatives must be handled.

\code{Element} is the alphabet of UI shapes. Each alternative is
one kind of thing the renderer can draw. The same exhaustiveness
property applies.

Together they form the program's two languages: one for
inputs, one for outputs. Both checked at compile time.

\subsection*{Why effects as data}

\code{Cmd} and \code{Sub} are descriptions of side effects,
not the effects themselves. The runtime interprets them. This
buys:

\begin{itemize}[leftmargin=*]
  \item Tests without the OS.
  \item Replay of any session.
  \item Composition without callbacks.
  \item A single place to add a new effect type.
\end{itemize}

\section{Sentence 2: the renderer}

\begin{quote}
The renderer is a six-stage pipeline (view, layout, paint,
diff, serialize, write) over a packed-cell canvas, fingerprinted
to skip whole frames when nothing changed.
\end{quote}

Part II touched the renderer at the edges; Part III walks
inside it. The architectural point: each stage is a pure
function from one value to another, the values are inspectable,
and the whole pipeline costs microseconds.

\subsection*{Why six stages, not one}

The same reason \code{init} and \code{update} are separate. Each
stage has a clear contract; bugs and optimisations are local.
Testing, profiling, replacement all happen one stage at a time.

\subsection*{Why packed cells}

A cell is a struct: codepoint + style id + flags. Eight or
sixteen bytes. The canvas is a flat array. Diff is SIMD over
the array. There is no widget hierarchy; only cells.

\subsection*{Why fingerprints}

A 64-bit hash of the element tree. Identical trees produce
identical hashes; the renderer skips the entire pipeline. An
idle app burns microseconds, not milliseconds.

\section{Sentence 3: the type system}

\begin{quote}
The DSL puts structural shape into the type system (via NTTPs
and concepts) and leaves content to runtime, so the compiler
catches shape errors and the cost of correctness is paid at
build time.
\end{quote}

Part II's longest chapter (the DSL) explored this. The
guarantees chapter listed ten of them. The C++26 chapter
explained how the language makes it possible.

\subsection*{Why structural in types}

The shape of a box (has-border, has-padding, has-width) does
not change at runtime. Encoding it in the type lets the
compiler check it. A box without a border that asks for a
border colour fails to build, not silently misrender at
runtime.

\subsection*{Why content at runtime}

Text strings, list lengths, model values: these depend on
input. Forcing them into types would balloon the type system
without buying safety. The DSL leaves these at runtime
(\code{text(s)}, runtime children) and only enforces shape.

\subsection*{Why both}

The DSL is the compile-time spelling; builders and runtime
constructors are the runtime spelling. They produce the same
\code{Element}. You pick by what your data looks like.

\section{What you should be able to do}

After Part II, you should:

\begin{itemize}[leftmargin=*]
  \item Write the four pure functions for a small app from
    scratch.
  \item Read \code{maya/include/maya/app/app.hpp} and
    understand it.
  \item Decode a typical DSL compile error and fix it.
  \item Pick the right pattern for a UI shape (master-detail,
    wizard, palette, log stream).
  \item Test \code{update} without a terminal.
  \item Reason about performance without profiling first.
\end{itemize}

If any of these is unclear, the relevant chapter is the place
to revisit.

\section{What Part III adds}

Part III opens the rendering pipeline. You will see:

\begin{itemize}[leftmargin=*]
  \item How a tree becomes a canvas (Chapter 14).
  \item How SIMD makes diffing sub-microsecond (Chapter 15).
  \item How ANSI bytes are serialised efficiently (Chapter 16).
  \item How inline mode coexists with the terminal scrollback
    (Chapter 17).
\end{itemize}

The architecture stays the same; Part III is the
implementation tour.

\section{What Part IV adds}

Part IV is reactivity and widgets:

\begin{itemize}[leftmargin=*]
  \item The signal/computed/effect system in depth.
  \item Events, focus, input parsing under the hood.
  \item \code{Cmd} and \code{Sub} reference.
  \item The widget catalogue.
  \item Building, packaging, shipping.
\end{itemize}

By the end of Part IV you will know enough to ship a non-
trivial \maya{} app in production.

\section{Where to take the book next}

Three suggestions:

\begin{itemize}[leftmargin=*]
  \item Build the task manager from Chapter~\ref{ch:end-to-end}.
    Force every concept to earn its place.
  \item Read \file{maya/include/maya/app/app.hpp} from top to
    bottom. The comments in the source extend the book.
  \item Ship something. Even a small tool. The architecture
    only feels right after the first 1000-line app you
    write.
\end{itemize}

The book is a guide; the apps you write are the test of
whether the guide worked.

\section*{Workshop}

\begin{enumerate}[leftmargin=*]
  \item Restate the three sentences in your own words. Compare
    with someone who has read the book.
  \item For each sentence, find a counter-example in another
    framework you have used.
  \item Imagine extending \maya{} with a new \code{Cmd}
    alternative. List every place you would need to change.
\end{enumerate}

\begin{recap}
Three sentences. Four functions, two sum types, two effect
languages, one pipeline, one canvas, types for shape and
runtime for content. Everything else in Part II is a
specialisation, a worked example, or a pattern. Internalising
the three sentences is the goal; the rest follows.
\end{recap}
